\section{基于预解式的谱距离估计}
\label{sec:error}

第~\ref{sec:numerics}~节的有限矩阵函数只能定位候选并生成一个分片重构场。小矩阵残量
本身不是连续算子的残量，更不能直接证明物理谱点的存在。本节固定同一候选、同一分片场及其全部
材料和几何输入，直接评价该场在真实界面上的连续迹跳跃；随后用正移位预解式的全界面迹
估计，把这些跳跃转化为到 $\sigma(A(\beta))$ 的距离上界。这个论证不先把分片场提升到
全局 $H^1$ 空间，也不要求数值求解任何预解式方程。

\subsection{重构场与界面跳跃}

固定第~\ref{sec:numerics}~节得到的 $\widehat k>0$，并在本节始终使用
\begin{align}
  \mu&:=\widehat k^2.
  \label{eq:resolvent-candidate-mu}
\end{align}
在 $\widehat k$ 处固定第~\ref{sec:numerics}~节候选计算交付的中心系数 $\xi$、左右
Cauchy 坐标 $c_-$、$c_+$、单胞边界密度、端口截断、连续核评价规则和完整单胞传播关系，
以及由这些数据确定的分片连续核场 $u_{h,M}$。以下称 $u_{h,M}$ 为原始场。它不是重新
离散后的场，也不随求积精度的提高而改变。近核向量只确定到非零复数倍，
因此先固定第~\ref{sec:numerics}~节使用的非零归一化与相位约定；后面的 $U$、界面缺陷和
全部单侧界都必须由这个同一尺度的场计算。

这种固定数据的方式适用于一般中心缺陷，并不限定密度的具体存储方式。设候选计算在均匀参数网格上
交付偶数个密度样本 $\{\eta_r\}_{r=0}^{N_\eta-1}$，其中 $N_\eta$ 是密度样本数。
固定这些样本，并按其存储位型将其解释为精确输入，再由数学离散 Fourier 变换定义有限三角密度：
\begin{align}
  \eta_m
  &:=\frac1{N_\eta}\sum_{r=0}^{N_\eta-1}\eta_r
    \exp\!\left(-\frac{2\pi\ii mr}{N_\eta}\right),
  \qquad
  -\frac{N_\eta}{2}\leq m\leq\frac{N_\eta}{2}-1.
  \label{eq:exact-density-dft}
\end{align}
该有限密度作用于使用同一个精确 $\mu$ 的完整数学准周期 Green 核，并与固定的有限
Rayleigh 和共同定义 $u_{h,M}$。可靠上界计算中的核展开截断、局部级数截断和曲线求积
只用于评价这个既定场的实际连续迹跳跃；它们的截断误差与舍入误差进入相应球上界，而不
重新定义另一个原始场。数值 FFT 也只可作为式~\eqref{eq:exact-density-dft} 的经验证实现，
不能独立定义密度系数。

原始场在每个均匀材料子区域内由固定密度和指定的精确 Helmholtz 核定义，并满足
\begin{align}
  (-\Delta-\mu\rho)u_{h,M}&=0
  \qquad\text{逐片成立}.
  \label{eq:raw-piecewise-pde}
\end{align}
它也精确满足 $y$ 方向的 $\beta$-准周期条件。本文假定
$u_{h,M}\in\mathcal H$，且它在每个分片直至相应界面局部属于 $H^2$。这些条件使
逐片 Green 恒等式和普通法向迹有意义，但不假定
$u_{h,M}\in H^1_\beta(B)$：有限边界积分匹配和端口截断仍可能在界面上留下小的
值跳跃和法向导数跳跃。

为避免重复计数，定义互不相交的界面族
\begin{align}
  \mathscr E&=\mathscr E_{\mathrm{mat}}\mathbin{\dot\cup}
  \mathscr E_{\mathrm{art}}.
  \label{eq:certification-interface-family}
\end{align}
$\mathscr E_{\mathrm{mat}}$ 包含 $\partial\Omega_0$（若
$\Omega_0\ne\varnothing$）以及 $\partial\Omega_-$、$\partial\Omega_+$ 的全部半无限周期
平移；$\mathscr E_{\mathrm{art}}$ 包含中心与首个周期胞元之间的 $\Gamma_-$、$\Gamma_+$，
以及左右半波导中所有相邻胞元的共享竖直界面。每条物理界面只出现一次。由于原始场的
准周期值与 $\partial_y$ 迹精确匹配，条带上下边界的 Green 项相消，不另计入
$\mathscr E$。

在每条 $\Gamma\in\mathscr E$ 上固定从侧 $1$ 指向侧 $2$ 的法向 $\nu_\Gamma$，并定义
\begin{align}
  a_\Gamma
  &:=u_{h,M}^{(2)}|_\Gamma
    -u_{h,M}^{(1)}|_\Gamma,
  &j_\Gamma
  &:=\partial_{\nu_\Gamma}u_{h,M}^{(1)}|_\Gamma
    -\partial_{\nu_\Gamma}u_{h,M}^{(2)}|_\Gamma.
  \label{eq:raw-interface-jumps}
\end{align}
在材料界面上，侧 $1$ 是介质柱内部、侧 $2$ 是外部，$\nu_\Gamma$ 由内指向外；因此
$a_\Gamma=u^{\mathrm{out}}-u^{\mathrm{in}}$，而
$j_\Gamma=\partial_{\nu_\Gamma}u^{\mathrm{in}}
-\partial_{\nu_\Gamma}u^{\mathrm{out}}$。在人工界面上，侧 $1$ 是较靠近中心的胞元，侧
$2$ 是较远的胞元，法向朝相应无穷远；在 $\Gamma_\pm$ 上这正是
$\nu_-=-e_x$、$\nu_+=e_x$。若改用共享界面两侧各自的外法向，$j_\Gamma$ 等于两侧
外法向导数之和。这个定向使后面的 Green 恒等式中值跳跃项和法向跳跃项都带正号。

选定参考长度 $\ell_*>0$，定义实际界面缺陷
\begin{align}
  \varepsilon_{\partial,\ell_*}^2
  &:={}
  \sum_{\Gamma\in\mathscr E}
  \left(
    \lVert j_\Gamma\rVert_{L^2(\Gamma)}^2
    +\ell_*^{-2}\lVert a_\Gamma\rVert_{L^2(\Gamma)}^2
  \right),
  \label{eq:boundary-defect-norm}\\
  U&:=\lVert u_{h,M}\rVert_{\mathcal H}.
  \label{eq:raw-field-mass}
\end{align}
两者使用同一个 $\mu$、同一个原始场和同一套材料参数。式
~\eqref{eq:boundary-defect-norm} 评价的是连续迹跳跃，不是 $J(\widehat k)\xi$、
$A(\widehat k)\xi$ 或端口矩阵不匹配量的范数。

无限半波导的求和仍保留完整传播矩阵。令 $P_\pm$ 表示把对应衰减 Cauchy 坐标向无穷远
推进一个周期的矩阵，则第 $j$ 个周期胞元的状态是 $P_\pm^jc_\pm$。一个胞元中的值和
法向跳跃作用连同 $L^2$ 权重形成正半定矩阵
$\mathsf G_{\partial,\pm}$ 后，半波导部分具有
\begin{align}
  \sum_{\varsigma\in\{-,+\}}
  \sum_{j=0}^{\infty}
  (P_\varsigma^jc_\varsigma)^*
  \mathsf G_{\partial,\varsigma}
  (P_\varsigma^jc_\varsigma)
  \label{eq:full-propagation-boundary-sum}
\end{align}
的结构；中心与首胞之间不能由平移生成的两条界面另行计算。每条共享界面只能归属于这两类
中的一类。式~\eqref{eq:full-propagation-boundary-sum} 使用完整 $P_\pm$，不把传播方向假定
为正交标量模态，也不删除非正规或 Jordan 耦合。材料跳跃可以直接在真实曲线上先作有符号
组合再求 $L^2$ 范数；下文所谓“非 Fourier 材料评价”仅指这一做法。端口场、人工竖直
界面以及源密度的构造仍然使用第~\ref{sec:numerics}~节的 Rayleigh/Fourier 表示。

\subsection{正移位预解式的全界面迹估计}

下面给出将式~\eqref{eq:boundary-defect-norm} 与全谱距离连接起来的几何常数。假定
$\mathscr E$ 局部有限，每条材料界面是简单闭合 $C^2$ 曲线，人工界面是准周期圆柱上的
竖直闭曲线。存在共同的法向邻域半宽 $\delta_{\mathrm{col}}>0$、曲率上界
$\kappa_*$ 和有限覆盖重数 $N_{\mathrm{ov}}$，使对每条界面的弧长参数
$\boldsymbol r_\Gamma(s)$，映射
\begin{align}
  \boldsymbol\Phi_\Gamma(s,t)
  &:=\boldsymbol r_\Gamma(s)+t\nu_\Gamma(s),
  &|t|&<\delta_{\mathrm{col}},
  \label{eq:interface-normal-neighborhood}
\end{align}
在该邻域内单射，并且任一点至多被 $N_{\mathrm{ov}}$ 个此类邻域覆盖。再要求
\begin{align}
  \vartheta_{\mathrm{col}}
  &:=1-\delta_{\mathrm{col}}\kappa_*>0,
  &c_0&:=\frac{N_{\mathrm{ov}}}
    {\vartheta_{\mathrm{col}}\delta_{\mathrm{col}}},
  &c_1&:=\frac{N_{\mathrm{ov}}\delta_{\mathrm{col}}}
    {2\vartheta_{\mathrm{col}}}.
  \label{eq:trace-geometry-constants}
\end{align}
这里 $\delta_{\mathrm{col}}$ 是几何法向邻域的半宽，不是第~\ref{sec:numerics}~节的
Nystr\"om 网格步长 $h$；$\vartheta_{\mathrm{col}}$ 是法向坐标 Jacobian 的下界，也不同于
式~\eqref{eq:background-bloch-fiber} 的 $x$ 向 Bloch 参数 $\theta$。单连通性本身并不保证
这些条件，计算可靠上界时还必须验证邻域单射性、不同界面及周期副本的分离性和覆盖重数界。

\begin{proposition}[全界面正移位迹估计]
  \label{prop:positive-shift-trace-bound}
  在上述几何条件以及 $0<\rho_-\leq\rho\leq\rho_+$ 下，对每个
  $v\in D(A(\beta))$，有
  \begin{align}
    T_{\ell_*}(v)^2
    &:={}
    \sum_{\Gamma\in\mathscr E}
    \left(
      \lVert v|_\Gamma\rVert_{L^2(\Gamma)}^2
      +\ell_*^2
       \lVert\partial_{\nu_\Gamma}v|_\Gamma\rVert_{L^2(\Gamma)}^2
    \right)
    \notag\\
    &\leq C_{\mathrm{tr},\ell_*}^2
      \lVert(A(\beta)+\mu I)v\rVert_{\mathcal H}^2,
    \label{eq:positive-shift-trace-bound}\\
    C_{\mathrm{tr},\ell_*}^2
    &:={}
    \frac{c_0}{\rho_-\mu^2}
    +\frac{c_1+\ell_*^2c_0}{4\mu}
    +\ell_*^2c_1\rho_+.
    \label{eq:trace-constant}
  \end{align}
\end{proposition}

\begin{proof}
  先说明所需的全局二阶控制。把 $v\in D(A(\beta))$ 以相位
  $\ee^{\ii\beta d}$ 沿 $y$ 延拓，并在 $y$ 方向用
  $\psi_m$、在 $x$ 方向用 Fourier 变换。由于 $v\in H^1(B)$ 且其分布 Laplacian 属于
  $L^2(B)$，Parseval 等式给出
  \begin{align*}
    \sum_m\int_{\RR}
    (\xi^2+\beta_m^2)^2|\widehat v_m(\xi)|^2\,\dd\xi<\infty.
  \end{align*}
  因而 $v\in H^2(B)$，并且
  $\lVert D^2v\rVert_{L^2(B)}=\lVert\Delta v\rVert_{L^2(B)}$。这一结论使用全局
  分布 Laplacian，故材料界面上的传输条件已经包含在 $D(A(\beta))$ 中。

  令 $f=(A(\beta)+\mu I)v$。由 $A(\beta)$ 的非负自伴性和谱定理，
  \begin{align*}
    \lVert v\rVert_{L^2(B)}^2
    &\leq\frac{1}{\rho_-\mu^2}\lVert f\rVert_{\mathcal H}^2,
    &\lVert\nabla v\rVert_{L^2(B)}^2
    &\leq\frac{1}{4\mu}\lVert f\rVert_{\mathcal H}^2,\\
    \lVert D^2v\rVert_{L^2(B)}^2
    &\leq\rho_+\lVert f\rVert_{\mathcal H}^2.
  \end{align*}
  最后一式来自 $\Delta v=-\rho A(\beta)v$、
  $\lVert A(\beta)v\rVert_{\mathcal H}\leq\lVert f\rVert_{\mathcal H}$ 及上述 Hessian
  恒等式。

  对 $g\in H^1(-\delta_{\mathrm{col}},\delta_{\mathrm{col}})$，一维基本定理和
  Cauchy--Schwarz 不等式给出
  \begin{align*}
    |g(0)|^2
    \leq\frac{1}{\delta_{\mathrm{col}}}
      \int_{-\delta_{\mathrm{col}}}^{\delta_{\mathrm{col}}}|g(t)|^2\,\dd t
    +\frac{\delta_{\mathrm{col}}}{2}
      \int_{-\delta_{\mathrm{col}}}^{\delta_{\mathrm{col}}}|g'(t)|^2\,\dd t.
  \end{align*}
  沿式~\eqref{eq:interface-normal-neighborhood} 的法线使用此式，利用 Jacobian 下界
  $\vartheta_{\mathrm{col}}$，再对任意有限界面子族求和并使用覆盖重数，得到
  \begin{align*}
    \sum_\Gamma\lVert v|_\Gamma\rVert_2^2
    &\leq c_0\lVert v\rVert_2^2+c_1\lVert\nabla v\rVert_2^2,\\
    \sum_\Gamma\lVert\partial_{\nu_\Gamma}v|_\Gamma\rVert_2^2
    &\leq c_0\lVert\nabla v\rVert_2^2+c_1\lVert D^2v\rVert_2^2.
  \end{align*}
  对递增有限子族取极限并代入前三个体范数估计，恰得
  式~\eqref{eq:positive-shift-trace-bound}--\eqref{eq:trace-constant}。
\end{proof}

后文简记
\begin{align}
  C&:=C_{\mathrm{tr},\ell_*},
  &Q&:=C\varepsilon_{\partial,\ell_*}.
  \label{eq:trace-defect-product}
\end{align}
$C$ 只依赖式~\eqref{eq:trace-constant} 中同一个 $\mu$、材料上下界、参考长度和经过验证的
几何输入；它不是从候选矩阵残量拟合出来的量。

\subsection{谱距离定理与可计算的可靠上界}

\begin{theorem}[原始界面缺陷的谱距离界]
  \label{thm:resolvent-spectral-distance}
  设 $A(\beta)$ 是式~\eqref{eq:fixed-beta-operator} 的非负自伴算子，界面族满足
  命题~\ref{prop:positive-shift-trace-bound} 的几何条件。设非零原始场
  $u_{h,M}\in\mathcal H$ 逐片属于 $H^2$，满足
  式~\eqref{eq:raw-piecewise-pde}、精确的 $\beta$-准周期条件，并且
  式~\eqref{eq:boundary-defect-norm} 有限。若
  \begin{align}
    Q&<U,
    \label{eq:positive-resolvent-denominator}
  \end{align}
  则
  \begin{align}
    \operatorname{dist}\!\left(\mu,\sigma(A(\beta))\right)
    &\leq\frac{2\mu Q}{U-Q}.
    \label{eq:resolvent-spectral-distance-bound}
  \end{align}
  这个结论是到整个算子谱的距离上界，不要求预先知道本质谱隙。
\end{theorem}

\begin{proof}
  若 $\mu\in\sigma(A(\beta))$，左端为零，结论成立。以下设
  $\mu\notin\sigma(A(\beta))$，并只为证明定义辅助场
  \begin{align}
    z&:=(A(\beta)-\mu I)^{-1}u_{h,M}
    \in D(A(\beta)).
    \label{eq:auxiliary-resolvent-field}
  \end{align}
  谱距离估计的计算不需求解 $z$。

  先证明所需的 Green 恒等式。在每个材料分片上对
  $u_{h,M}$ 和 $z$ 使用第二 Green 恒等式。共享界面两侧的外法向分别为
  $\nu_\Gamma$ 和 $-\nu_\Gamma$，而 $z\in D(A(\beta))\subset H^2(B)$ 的值迹和普通
  法向导数迹均单值。结合式~\eqref{eq:raw-interface-jumps}，两侧边界项之和恰为
  \begin{align*}
    j_\Gamma\overline z
    +a_\Gamma\partial_{\nu_\Gamma}\overline z.
  \end{align*}
  为处理无穷远，取实值截断函数 $\chi_R(x)$，使其在 $|x|\leq R$ 时为一、在
  $|x|\geq2R$ 时为零，并满足
  $\lVert\chi_R'\rVert_\infty=O(R^{-1})$、
  $\lVert\chi_R''\rVert_\infty=O(R^{-2})$。以 $\chi_Rz$ 作第二个场并先对有限个分片求和。
  顶、底边界项由精确准周期条件相消；含 $\chi_R'$、$\chi_R''$ 的体项由
  $u_{h,M}\in\mathcal H$ 和 $z\in H^2(B)$ 控制并趋于零，界面上的附加截断项则由
  $\sum_\Gamma\lVert a_\Gamma\rVert_2^2<\infty$ 与命题
  ~\ref{prop:positive-shift-trace-bound} 控制并趋于零。主界面级数由
  式~\eqref{eq:boundary-defect-norm}、该命题和 Cauchy--Schwarz 不等式绝对收敛。因此令
  $R\to\infty$ 得到
  \begin{align}
    U^2
    =\sum_{\Gamma\in\mathscr E}
    \int_\Gamma
    \left(
      j_\Gamma\overline z
      +a_\Gamma\partial_{\nu_\Gamma}\overline z
    \right)\,\dd s.
    \label{eq:raw-green-identity}
  \end{align}

  记
  $\delta=\operatorname{dist}(\mu,\sigma(A(\beta)))>0$。由自伴算子的谱定理，
  $\lVert(A(\beta)-\mu I)^{-1}\rVert=\delta^{-1}$。又由
  $(A(\beta)+\mu I)z=u_{h,M}+2\mu z$，
  \begin{align}
    \lVert(A(\beta)+\mu I)z\rVert_{\mathcal H}
    &\leq\left(1+\frac{2\mu}{\delta}\right)U.
    \label{eq:positive-shift-of-resolvent}
  \end{align}
  对式~\eqref{eq:raw-green-identity} 使用带 $\ell_*$ 权重的 Cauchy--Schwarz 不等式和
  命题~\ref{prop:positive-shift-trace-bound}，得到
  \begin{align*}
    U^2
    &\leq\varepsilon_{\partial,\ell_*}T_{\ell_*}(z)
    \leq Q\left(1+\frac{2\mu}{\delta}\right)U.
  \end{align*}
  除以 $U>0$，再利用 $U-Q>0$ 移项，即有
  $\delta(U-Q)\leq2\mu Q$，这就是式
  ~\eqref{eq:resolvent-spectral-distance-bound}。
\end{proof}

用数值近似评价上述理论公式可得到谱距离估计；若要得到可靠上界，则必须为定理输入建立
可靠单侧界。具体地，设
\begin{align}
  C&\leq\overline C,
  &\varepsilon_{\partial,\ell_*}&\leq\overline\varepsilon_\partial,
  &0<\underline U&\leq U,
  \label{eq:certified-one-sided-inputs}\\
  \overline Q&:=\overline C\,\overline\varepsilon_\partial,
  &\underline D&:=\underline U-\overline Q>0.
  \label{eq:certified-positive-denominator}
\end{align}
$\overline C$ 来自材料、候选 $\mu$ 和法向邻域几何的向上包围；
$\overline\varepsilon_\partial$ 来自同一原始场全部实际值跳跃和法向跳跃的向上包围；
$\underline U$ 则是该场加权质量的向下包围。三者不得混用不同候选、不同归一化或重新求解
后的密度。由定理~\ref{thm:resolvent-spectral-distance} 立即得到可计算的可靠上界
\begin{align}
  R_{\mathrm{spec}}
  &:=\frac{2\mu\overline C\,\overline\varepsilon_\partial}
    {\underline U-\overline C\,\overline\varepsilon_\partial},
  &\operatorname{dist}\!\left(\mu,\sigma(A(\beta))\right)
  &\leq R_{\mathrm{spec}}.
  \label{eq:certified-spectral-distance}
\end{align}
正分母条件~\eqref{eq:certified-positive-denominator} 是有限界成立的必要检查；若它不成立，
公式不提供有限谱距离界。$R_{\mathrm{spec}}$ 是到算子谱的距离上界，不是残量范数。
因此向外舍入的区间
\begin{align}
  I_{\mathrm{spec}}
  &:=[\mu-R_{\mathrm{spec}},\mu+R_{\mathrm{spec}}]
  \label{eq:certified-spectral-interval}
\end{align}
必与 $\sigma(A(\beta))$ 相交。这个结论既不识别最近谱点的编号或重数，也不保证区间内只有
一个谱点。

若另行证明某个开区间 $G$ 与 $\sigma_{\mathrm{ess}}(A(\beta))$ 不交，并且
$I_{\mathrm{spec}}\subset G$，则
命题~\ref{prop:discrete-spectrum-in-gap} 才允许把其中至少一个谱点解释为孤立、有限重的
导模特征值。本文的谱距离可靠上界本身不证明这样的 $G$ 存在，也不把猜想
~\ref{conj:essential-spectrum} 当成数值带隙的证明。

\paragraph{体缺陷的独立扩展。}
当前连续核原始场逐片精确满足式~\eqref{eq:raw-piecewise-pde}，故体缺陷为零。若将来使用
不精确满足 PDE 的分片场，令
$r=(-\Delta-\mu\rho)u_{h,M}$ 且
$\varepsilon_{\mathrm{vol}}:=\lVert\rho^{-1}r\rVert_{\mathcal H}$，则同一证明在
式~\eqref{eq:raw-green-identity} 中增加体积分，并给出
\begin{align*}
  \operatorname{dist}\!\left(\mu,\sigma(A(\beta))\right)
  \leq
  \frac{2\mu\overline C\,\overline\varepsilon_\partial
    +\overline\varepsilon_{\mathrm{vol}}}
  {\underline U-\overline C\,\overline\varepsilon_\partial}.
\end{align*}
这一扩展不参与本文算例，也不改变精确 PDE 主线。
